\begin{lemma}
\label{lemma-homotopy-equivalence}
Let $f : X \to Y$ be a map of simplicial abelian groups. If $f$ is a
homotopy equivalence of simplicial sets, then $f$ induces a
quasi-isomorphism of associated chain complexes.
\end{lemma}

\begin{proof}
In this proof we will write $H_n(Z) = H_n(s(Z)) = H_n(N(Z))$
when $Z$ is a simplicial abelian group, with $s$ and $N$ as in
Section \ref{section-complexes}.
Let $\mathbf{Z}[X]$ denote the free abelian group on $X$ viewed as
a simplicial set and similarly for $\mathbf{Z}[Y]$. Consider the
commutative diagram
$$
\xymatrix{
\mathbf{Z}[X] \ar[r]_g \ar[d] &
\mathbf{Z}[Y] \ar[d] \\
X \ar[r]^f & Y
}
$$
of simplicial abelian groups. Since taking the free abelian group on a set
is a functor, we see that the horizontal arrow is a homotopy equivalence of
simplicial abelian groups, see Lemma \ref{lemma-functorial-homotopy}.
By Lemma \ref{lemma-homotopy-equivalence-s-N}
we see that $H_n(g) : H_n(\mathbf{Z}[X]) \to H_n(\mathbf{Z}[Y])$ is bijective
for all $n \geq 0$.

\medskip\noindent
Let $\xi \in H_n(Y)$. By definition of $N(Y)$ we can represent
$\xi$ by an element $y \in N(Y_n)$ whose boundary is zero.
This means $y \in Y_n$ with $d^n_0(y) = \ldots = d^n_{n - 1}(y) = 0$
because $y \in N(Y_n)$ and $d^n_n(y) = 0$ because the boundary
of $y$ is zero. Denote $0_n \in Y_n$ the zero element. Then we see that
$$
\tilde y = [y] - [0_n] \in (\mathbf{Z}[Y])_n
$$
is an element with $d^n_0(\tilde y) = \ldots = d^n_{n - 1}(\tilde y) = 0$
and $d^n_n(\tilde y) = 0$. Thus $\tilde y$ is in $N(\mathbf{Z}[Y])_n$
has boundary $0$, i.e., $\tilde y$ determines a
class $\tilde \xi \in H_n(\mathbf{Z}[Y])$ mapping to $\xi$.
Because $H_n(\mathbf{Z}[X]) \to H_n(\mathbf{Z}[Y])$ is bijective
we can lift $\tilde \xi$ to a class in $H_n(\mathbf{Z}[X])$.
Looking at the commutative diagram above
we see that $\xi$ is in the image of $H_n(X) \to H_n(Y)$.

\medskip\noindent
Let $\xi \in H_n(X)$ be an element mapping to zero in $H_n(Y)$.
Exactly as in the previous parapgraph we can represent
$\xi$ by an element $x \in N(X_n)$ whose boundary is zero, i.e.,
$d^n_0(x) = \ldots = d^n_{n - 1}(x) = d^n_n(x) = 0$.
In particular, we see that $[x] - [0_n]$ is an element of
$N(\mathbf{Z}[X])_n$ whose boundary is zero, whence defines
a lift $\tilde \xi \in H_n(\mathbf{Z}[x])$ of $\xi$.
The fact that $\xi$ maps to zero in $H_n(Y)$
means there exists a $y \in N(Y_{n + 1})$ whose boundary is $f_n(x)$.
This means $d^{n + 1}_0(y) = \ldots = d^{n + 1}_n(y) = 0$ and
$d^{n + 1}_{n + 1}(y) = f(x)$.
However, this means exactly that
$z = [y] - [0_{n + 1}]$ is in $N(\mathbf{Z}[y])_{n + 1}$ and
$$
g([x] - [0_n]) = [f(x)] - [0_n] = \text{boundary of }z
$$
This proves that $\tilde \xi$ maps to zero in $H_n(\mathbf{Z}[y])$.
As $H_n(\mathbf{Z}[X]) \to H_n(\mathbf{Z}[Y])$ is bijective
we conclude $\tilde \xi = 0$ and hence $\xi = 0$.
\end{proof}